% Curated cumulative working-reader wrapper through HI-OLP-0009.
% Source authority: OpenLogicProject/OpenLogic commit
% 9620cc73f9c8e0ad003c514a5d3748f29611c4c0.
% This wrapper adds only Hindi title/status/navigation matter; each translated
% section remains an independently hashed source-preserving subfile.

\documentclass[../../../include/open-logic-section]{subfiles}
\usepackage{flafter}
\usepackage{float}

\renewcommand*{\ifgitinfo}{\sffamily\footnotesize संशोधन~9620cc7}
\renewcommand*{\photocredits}{}
\renewcommand*{\contentsname}{विषय-सूची}
\floatplacement{figure}{H}

\hypersetup{
  pdftitle={ओपन लॉजिक परियोजना: हिन्दी कार्यशील अनुवाद},
  pdfauthor={Interlanguage Hindi translation programme},
  pdfsubject={समुच्चयों और संबंधों पर नौ अनूदित खंड},
  pdfkeywords={तर्कशास्त्र, समुच्चय, संबंध, हिन्दी, Open Logic Project}
}

\begin{document}

\begingroup
\thispagestyle{empty}
\centering
\vspace*{0.13\textheight}
{\Huge\bfseries ओपन लॉजिक परियोजना\par}
\vspace{0.8cm}
{\LARGE हिन्दी कार्यशील अनुवाद\par}
\vspace{1.2cm}
{\Large समुच्चय और संबंध\par}
\vspace{0.4cm}
{\large नौ आरम्भिक खंड\par}
\vfill
{\normalsize स्रोत: Open Logic Project, संशोधन
\texttt{9620cc73f9c8e0ad003c514a5d3748f29611c4c0}\par}
\vspace{0.5cm}
{\small यह मशीन-सहायित कार्यशील अनुवाद है। प्रत्येक सम्मिलित खंड का
गणितीय विन्यास जाँचा गया है, PDF सफलतापूर्वक बनाया गया है और प्रत्येक
पृष्ठ का दृश्य परीक्षण किया गया है। यह पूर्ण पुस्तक, मानवीय समीक्षा,
आलोचनात्मक संस्करण अथवा PDF/UA सुगम्यता का दावा नहीं है। अगला अनुवाद
खंड \texttt{orders.tex} है।\par}
\vspace*{0.08\textheight}
\par
\endgroup
\clearpage
\raggedbottom

\tableofcontents*
\clearpage

\subfile{basics}
\subfile{subsets}
\subfile{important-sets}
\subfile{unions-and-intersections}
\subfile{pairs-and-products}
\subfile{russells-paradox}
\subfile{../relations/relations-as-sets}
\subfile{../relations/special-properties}
\subfile{../relations/equivalence-relations}

\end{document}
